% Generated draft: requires mathematical and source review.
\documentclass[11pt]{article}
\usepackage[margin=25mm]{geometry}
\usepackage{fontspec}
\setmainfont{Latin Modern Roman}
\usepackage{amsmath,amssymb,mathtools,unicode-math}
\setmathfont{Latin Modern Math}
\usepackage{enumitem,needspace,xurl,hyperref}
\setlength{\parindent}{0pt}
\setlength{\parskip}{5pt}
\setlength{\emergencystretch}{3em}
\newcommand{\N}{\mathbb{N}}
\newcommand{\Z}{\mathbb{Z}}
\newcommand{\Q}{\mathbb{Q}}
\newcommand{\R}{\mathbb{R}}
\newcommand{\C}{\mathbb{C}}
\newcommand{\F}{\mathbb{F}}
\newcommand{\A}{\mathbb{A}}
\newcommand{\PP}{\mathbb{P}}
\newcommand{\E}{\mathbb{E}}
\newcommand{\eps}{\varepsilon}
\newcommand{\dd}{\,\mathrm{d}}
\begin{document}
\section*{UnsolvedMath ID: 30003931}
\textbf{Source problem code:} catalogue entry
\subsection{Definitions and mathematical statement}
Learning-and-optimization (LEO), also termed learning enabled optimization in the supplied references, combines statistical learning and optimization. The supplied problem is a request for a conceptual and mathematical framework, rather than a conjecture about a specified class of models.

To explain the proposed probability criterion schematically, suppose a minimization model has a nonempty feasible set \(X\), an objective \(J:X\to\R\) with finite infimum \(J_*:=\inf_{x\in X}J(x)\), and a measurable random feasible output \(\widehat x\) on a probability space with law \(\mathbb P\). For a specified tolerance \(\varepsilon\geq 0\), an output is \(\varepsilon\)-optimal when \(J(\widehat x)\leq J_*+\varepsilon\); its probability of being \(\varepsilon\)-optimal is
\[
\mathbb P\!\left(J(\widehat x)\leq J_*+\varepsilon\right).
\]
An exact deterministic optimum would satisfy \(J(x)=J_*\). This notation is illustrative only: the source does not specify a minimization formulation, feasible set, objective, probability law, measurability assumptions, or finite optimal value. It also does not specify whether randomness comes from training data, an algorithm, future observations, or some combination, or whether optimality is measured for an empirical or population objective.

\textbf{Research problem.} For LEO models, determine whether deterministic optimality is the appropriate goal or whether an appropriate notion of statistical optimality should quantify the probability of obtaining an \(\varepsilon\)-optimal solution. Formulate a suitable notion of generalizability and determine whether hypothesis tests can validate such models. The source supplies no model class, generalization target, null and alternative hypotheses, sampling assumptions, confidence requirement, or error criterion. Consequently, these requests cannot yet be expressed as a uniquely specified quantified theorem or testing problem; choosing those ingredients is part of the requested framework.
\subsection{Short English statement}
Develop an appropriate notion of optimality for LEO models, possibly using the probability of obtaining an \(\varepsilon\)-optimal solution, and define generalizability and statistical model validation. The source poses a research program without enough assumptions to specify a unique theorem.
\subsection{Sources}
\begin{itemize}
\item Supplied local UnsolvedMath registry record, permanent ID \(30003931\), title: \emph{Statistical Optimality and Validation in Learning-and-Optimization Models}; source problem code: \texttt{catalogue entry}. The collection is \emph{Oberwolfach Reports Open Problems}. No individual proposer is supplied; the proposed year is recorded as 2018.
\item Supplied report citation: \emph{New Directions in Stochastic Optimisation} (2019), Oberwolfach Reports, pp.\ 2303--2384, DOI 10.4171/owr/2018/38. \url{https://doi.org/10.4171/owr/2018/38}. The supplied background describes these reports as documenting mathematical workshops and their open-problem discussions.
\item \emph{Learning Enabled Optimization: Towards a Fusion of Statistical Learning and Stochastic Optimization}. Supplied literature link: \url{https://optimization-online.org/2017/03/5904/}.
\item \emph{Distribution-free Algorithms for Learning Enabled Optimization}. Supplied literature link: \url{https://optimization-online.org/2020/03/7661/}.
\item Literature assessment dated exactly \texttt{2026-08-21}: the supplied curation labels the problem \texttt{partially\_solved}. It reports that the LEO literature proposes statistical optimality, generalization error, model-fidelity tests, and model-selection procedures, while a complete theory of generalization and validation for arbitrary LEO models remains open. This is supplied dated literature triage derived from the OWR curation, not an independently verified assessment; the time-sensitive status and cited claims require independent verification.
\item The registry labels the statement \texttt{corrected\_verified} and states that mathematical content was preserved while notation, typography, and omitted definitions were repaired for self-contained publication. These are supplied editorial assertions, not independent mathematical verification of this draft.
\end{itemize}
\end{document}
